\documentclass[../../main.tex]{subfiles}
\begin{document}

{\bf[解]}

$\sin t, \cos t$の周期性から
\begin{align}
  0\le t<2\pi \label{eq:1}
\end{align}
で考える．題意より$P(X,Y)$は
\begin{align}
  \begin{cases}
    X = \sin t + \cos t \\
    Y = k\sin^2t\cos^2t
  \end{cases} \label{eq:2}
\end{align}
である．$f(t)=|OP|^2$とすると
\begin{align}
  f(t)
  &=X^2+Y^2 \\
  &=(\sin t+\cos t)^2+(k\sin^2t \cos^2t)^2 \\
  &=k^2p^4+2p+1\equiv g(p) \label{eq:3}
\end{align}
である．ただし$p=\sin t\cos t$とした．
この時$p=\dfrac{1}{2}\sin\dfrac{t}{2}$だから\cref{eq:1}と合わせて
\begin{align}
  \dfrac{-1}{2}\le p\le\dfrac{1}{2} \label{eq:4}
\end{align}
を満たす．

$f(t)=g(p)$の挙動を調べるために一階微分を計算すると
\begin{align}
  g'(p)
  &=4k^2p^3+2 = 4k^2\left(p^2+\dfrac{1}{2k^2}\right) \\
  &=4k^2\left(p+\alpha\right)\left(p^2-p\alpha+\alpha^2\right) \label{eq:5}
\end{align}
である．ただし
\begin{align}
  \alpha=\sqrt[3]{\dfrac{1}{2k^2}} \label{eq:6}
\end{align}
とおいた．ここで$p^2-p\alpha+\alpha^2>0$だから，$g'(p)$の符号は$p+\alpha$の符号による．
従って\cref{eq:4}より$g(p)$の増減表は$\alpha$の値に応じて以下のようになる．

\subparagraph{(i)$\dfrac{1}{2}\ge \alpha$つまり$2\le k$の時}

この場合は$p+\alpha=0$となる$p$が存在して$g(p)$の増減表は以下のようになる．

\begin{table}[htb]
  \centering
  \caption{$g(p)$ の増減表}
  \label{tab:g_increase_decrease}
  \renewcommand{\arraystretch}{1.4}
  \begin{tabular}{c|ccccc}
    $p$ & $-\dfrac{1}{2}$ & $\cdots$ & $-\alpha$ & $\cdots$ & $\dfrac{1}{2}$ \\
    \hline
    $g'(p)$ & & $-$ & $0$ & $+$ & \\
    \hline
    $g(p)$ & $\dfrac{k^2}{16}$ & $\searrow$ & 極小 & $\nearrow$ & $\dfrac{k^2}{16} + 2$ \\
    &                   &           & $-\dfrac{3\alpha}{2} + 1$ & &
  \end{tabular}
\end{table}

従って最大最小値は
\begin{align}
  \begin{cases}
    \max g= g(1/2)     = \dfrac{k^2}{16}+2  \\
    \min g= g(-\alpha) = 1-\dfrac{3}{2}\sqrt[3]{\dfrac{1}{2k^2}}
  \end{cases}
\end{align}
となる．

\subparagraph{(ii)$\dfrac{1}{2}\le \alpha$つまり$2\ge k>0$の時}

この時は$p+\alpha \ge 0$より$g'(p)\ge0$であって$g(p)$は単調増加である．
従って最大最小値は
\begin{align}
  \begin{cases}
    \max g=g\left(\dfrac{1}{2}\right) = \dfrac{k^2}{16}+2 \\
    \min g=g\left(\dfrac{-1}{2}\right) = \dfrac{k^2}{16}
  \end{cases}
\end{align}
である．
\casesend

以上から求める最大小値は，
\begin{align}
  \begin{cases}
    \max g=\dfrac{k^2}{16}+2 &  \min g=\dfrac{k^2}{16} \quad (0<k\le 2\text{の時}) \\
    \max g=\dfrac{k^2}{16}+2 & \min g=1-\dfrac{3}{2}\sqrt[3]{\dfrac{1}{2k^2}} \quad (2\le k\text{の時})
  \end{cases}
\end{align}
である．
\newpage
\end{document}
